% P3CTeX optional GNT integration regression.
% Compile from tex/tests with TEXINPUTS="../latex:../code:".
\documentclass[PR,CAT,default,project-management]{P3CTeX}
\makeatletter
\@ifpackageloaded{pxUML}{}{\PackageError{pxPM-test}{default module not loaded}{}}
\@ifpackageloaded{pxPM}{}{\PackageError{pxPM-test}{GNT module not loaded}{}}
\makeatother
\PECTeXconfig{cover=false,toc=false,data={student={Estudiant},grade={Enginyeria Informàtica},subj-fullname={Projecte Gantt},ca-fullname={Planificació},ca-shortname={PAC1}}}
\begin{document}
\section{Pla de treball}
\pxGantt{start=01/10/2026,end=31/01/2027,days-per-block=35,day-width=6mm}
\pxGanttItem{Investigació}{01/10/2026}{20/10/2026}
\pxGanttItem{Disseny}{15/10/2026}{10/11/2026}
\pxGanttItem{Implementació}{01/11/2026}{15/12/2026}
\pxGanttItem{Avaluació de models locals}{05/12/2026}{15/01/2027}
\pxGanttItem{Memòria final}{01/01/2027}{31/01/2027}
\pxGanttEnd
\section{Continuació}
La pàgina manté la seva orientació original.
\end{document}
